\section{Evaluasi}
\label{sec:evaluasi-bab16}

\subsection{Soal Pilihan Ganda}

\begin{enumerate}
    \item Perintah pada MS-DEBUG yang dipakai untuk membongkar kode
          mesin menjadi \eng{mnemonic} adalah \dots
          \begin{enumerate}
              \item \code{T} \item \code{U}
              \item \code{D} \item \code{G}
          \end{enumerate}

    \item Nama variabel \code{larik} tidak muncul pada hasil
          \eng{disassembly}, melainkan berubah menjadi
          \lstasm|[SI+01F4]|. Penyebabnya adalah \dots
          \begin{enumerate}
              \item nama simbolik tidak disimpan di dalam berkas biner
              \item DEBUG selalu menghapus nama simbolik
              \item linker mengubah nama simbolik menjadi alamat
                    relatif
              \item disassembler hanya mampu membaca register
          \end{enumerate}

    \item Instruksi tunggal yang tepat menggantikan
          \lstasm|MUL BX| untuk mengalikan \reg{AX} (tak bertanda)
          dengan delapan adalah \dots
          \begin{enumerate}
              \item \lstasm|SHR AX, 3| \item \lstasm|SHL AX, 3|
              \item \lstasm|SAR AX, 3| \item \lstasm|DIV AX, 8|
          \end{enumerate}

    \item Ciri paling dapat diandalkan bahwa suatu blok pada hasil
          \eng{disassembly} adalah perulangan adalah \dots
          \begin{enumerate}
              \item adanya instruksi \asm{INT 21h}
              \item lompatan mundur ke alamat yang lebih kecil
              \item adanya \lstasm|MOV AX, @DATA|
              \item lompatan maju yang jauh
          \end{enumerate}

    \item Keunggulan \lstasm|XOR AX, AX| dibandingkan
          \lstasm|MOV AX, 0| untuk mengosongkan register adalah \dots
          \begin{enumerate}
              \item hasilnya selalu nol pada semua kondisi
              \item lebih pendek satu byte dan lebih cepat satu
                    siklus
              \item aman dipakai di antara \asm{CMP} dan lompatan
                    bersyarat
              \item sekaligus mengosongkan \reg{DX}
          \end{enumerate}
\end{enumerate}

\subsection{Soal Uraian}

\begin{enumerate}
    \item Jelaskan tiga keterbatasan \eng{disassembly}. Untuk
          masing-masing keterbatasan, berikan satu contoh potongan
          byte atau instruksi yang memperlihatkan mengapa
          \eng{disassembly} sulit ditafsirkan.

    \item Perhatikan potongan hasil \eng{disassembly} berikut.

\begin{lstlisting}[language=MASM, numbers=none]
1D2F:0100 B90A00    MOV  CX,000A
1D2F:0103 B80000    MOV  AX,0000
1D2F:0106 0304      ADD  AX,[SI]
1D2F:0108 83C602    ADD  SI,+02
1D2F:010B E2F9      LOOP 0106
1D2F:010D D1E0      SHL  AX,1
1D2F:010F D1E0      SHL  AX,1
1D2F:0111 D1E0      SHL  AX,1
1D2F:0113 8906F401  MOV  [01F4],AX
\end{lstlisting}

          \begin{enumerate}
              \item Sebutkan pola yang tampak dan tunjukkan buktinya.
              \item Hitung jumlah instruksi yang dieksekusi untuk
                    seluruh program tersebut.
              \item Hitung perkiraan jumlah siklusnya, lalu
                    sebutkan satu perubahan pada potongan itu yang
                    masih dapat menghemat siklus.
          \end{enumerate}

    \item Diberikan program naif berikut: akumulator di memori,
          pengali 16 dimuat ke \reg{BX} setiap iterasi, perkalian
          memakai \asm{MUL}, dan penghitung perulangan memakai
          \asm{DEC} diikuti \asm{JNZ}. Tuliskan versi teroptimasi
          lengkap untuk program tersebut, sebutkan nama setiap
          teknik yang Anda pakai, dan jelaskan syarat berlaku
          masing-masing. Sebutkan pula metrik yang Anda pakai untuk
          membuktikan perbaikannya.

    \item Seorang mahasiswa mengganti sebuah \asm{MUL} dengan
          pergeseran, lalu melaporkan bahwa programnya menjadi dua
          kali lebih cepat. Rancang langkah-langkah pengukuran yang
          harus ia lakukan agar klaim tersebut dapat dipercaya,
          kemudian jelaskan alasan mengapa optimasi dini justru
          lebih sering merugikan daripada menguntungkan.
\end{enumerate}

\subsection{Kisi-kisi Penilaian}

\begin{table}[H]
    \centering
    \caption{Kisi-kisi evaluasi Bab~\ref{chap:disassembly}.}
    \label{tab:kisi-bab16}
    \small
    \begin{tabular}{@{}p{1.5cm}p{5.0cm}p{2.6cm}c@{}}
        \toprule
        \textbf{Sub-CPMK} & \textbf{Indikator yang diukur} &
        \textbf{Bentuk soal} & \textbf{Bobot} \\
        \midrule
        4.3 & Membaca hasil \eng{disassembly} dan mengenali polanya &
        Pilihan ganda 1, 2, 4; Uraian 1, 2 & 30\% \\
        \addlinespace
        4.3 & Menjelaskan keterbatasan \eng{disassembly} &
        Pilihan ganda 2; Uraian 1 & 15\% \\
        \addlinespace
        4.3 & Menghitung metrik optimasi &
        Uraian 2 bagian b dan c & 20\% \\
        \addlinespace
        4.3 & Menerapkan teknik optimasi level instruksi &
        Pilihan ganda 3, 5; Uraian 3 & 25\% \\
        \addlinespace
        4.3 & Merancang pengukuran dan memahami batas optimasi &
        Uraian 4 & 10\% \\
        \bottomrule
    \end{tabular}
\end{table}

\begin{checklistbox}
\textbf{Checklist Pencapaian Kompetensi Bab~\ref{chap:disassembly}}

Seorang mahasiswa dinyatakan kompeten pada Sub-CPMK 4.3 apabila:
\begin{itemize}
    \item[\checked] dapat membongkar program dengan perintah
          \code{U} serta membaca alamat, byte mesin, dan
          \eng{mnemonic} dengan benar;
    \item[\checked] dapat menyebutkan keterbatasan \eng{disassembly}
          dan menjelaskan mengapa kode serta data sulit dibedakan;
    \item[\checked] dapat mengenali minimal lima pola hasil
          \eng{disassembly} beserta cirinya;
    \item[\checked] dapat menghitung empat metrik optimasi dan
          menjelaskan bahwa tanpa pengukuran optimasi hanyalah
          dugaan;
    \item[\checked] dapat menerapkan minimal tiga teknik optimasi
          level instruksi beserta syarat berlakunya;
    \item[\checked] dapat menjelaskan trade-off ukuran dan kecepatan
          serta memilih dengan alasan;
    \item[\checked] berhasil membuktikan perbaikan program melalui
          pengukuran yang adil.
\end{itemize}
\end{checklistbox}
